\documentclass[../../main.tex]{subfiles}
\begin{document}


{\bf [解]}
題意の二つの球の方程式は
\begin{align}
 \begin{cases}
  C_1: x^2+y^2+z^2=r^2 \\
  C_2: (x-1)^2+y^2+z^2=1-r^2 
 \end{cases} \label{eq:1}
\end{align}
である．辺々引くと
\begin{align}
 2x-1 = 2r^2-1 \quad x = r^2 \label{eq:2}
\end{align}
だからこれが交面である．

\begin{figure}[htb]
\begin{center}
\begin{tikzpicture}[scale=3.2, >=stealth]

  % --- 1. パラメータ設定 (r = 1/sqrt(2) ≒ 0.7071) ---
  \pgfmathsetmacro{\rVal}{1/sqrt(2)}
  \pgfmathsetmacro{\rSq}{\rVal*\rVal} % x = r^2 = 0.5
  \pgfmathsetmacro{\rTwo}{sqrt(1 - \rSq)} % 半径 sqrt(1-r^2) = 1/sqrt(2)
  \pgfmathsetmacro{\yIntersect}{sqrt(\rSq - \rSq*\rSq)} % y = sqrt(r^2 - r^4) = 0.5

  % 中心座標
  \coordinate (O1) at (0, 0);
  \coordinate (O2) at (1, 0);

  % --- 2. 共通部分の斜線ハッチング ---
  % 円 C1 の右側 (x >= r^2) かつ 円 C2 の左側 (x <= r^2)
  \begin{scope}
    \clip (O1) circle (\rVal);
    \clip (O2) circle (\rTwo);
    \fill[pattern=north east lines, pattern color=blue!70] 
      (-0.2, -0.9) rectangle (1.2, 0.9);
  \end{scope}

  % --- 3. 座標軸 ---
  \draw[->, thick] (-0.9, 0) -- (1.9, 0) node[right] {$x$};
  \draw[->, thick] (0, -0.9) -- (0, 0.9) node[above] {$y$};
  \node[below left, font=\small] at (O1) {$O$};

  % --- 4. 2つの円の描画 ---
  % 円 C1 (中心 原点, 半径 r)
  \draw[ultra thick, blue!80!black] (O1) circle (\rVal);
  \node[blue!80!black, font=\small, above left] at (-0.45, 0.5) {$C_1: x^2+y^2=r^2$};

  % 円 C2 (中心 (1,0), 半径 sqrt(1-r^2))
  \draw[ultra thick, red!80!black] (O2) circle (\rTwo);
  \node[red!80!black, font=\small, above right] at (1.45, 0.5) {$C_2: (x-1)^2+y^2=1-r^2$};

  % --- 5. 交線 x = r^2 (破線) ---
  \draw[thick, dashed, purple!80!black] (\rSq, -0.75) -- (\rSq, 0.75)
    node[above, font=\small] {交線 $x = r^2$};

  % --- 6. 各点のプロットと目盛り ---
  % 中心 O2(1, 0)
  \fill[black] (O2) circle (1.0pt) node[below] {$1$};

  % 交点の x 座標 x = r^2
  \fill[purple!80!black] (\rSq, 0) circle (1.0pt) node[below left, font=\small] {$r^2$};

  % 2つの交点
  \coordinate (IntTop) at (\rSq, \yIntersect);
  \coordinate (IntBot) at (\rSq, -\yIntersect);
  \fill[purple!80!black] (IntTop) circle (1.2pt);
  \fill[purple!80!black] (IntBot) circle (1.2pt);

  % 半径の注記 (C1: 半径 r, C2: 半径 sqrt(1-r^2))
  \draw[<->, thin, blue!80!black] (O1) -- ({-\rVal*cos(45)}, {-\rVal*sin(45)})
    node[midway, above left, font=\scriptsize] {$r$};
  \draw[<->, thin, red!80!black] (O2) -- ({1 + \rTwo*cos(45)}, {-\rTwo*sin(45)})
    node[midway, above right, font=\scriptsize] {$\sqrt{1-r^2}$};

\end{tikzpicture}
\end{center}
\caption{$xy$ 平面における2つの円 $C_1, C_2$ と共通部分および交線 $x=r^2$}
\end{figure}

\paragraph{(1)}

共通部分の立体は，球$C_1$の$x\ge r^2$の部分と，球$C_2$の$x\le r^2$の部分の和である．球$i$を単位球に正規化した時の対応する球欠の体積（$\pi$倍する前）を$V_i$とすると，スケーリングにより
\begin{align}
V(r)=r^3V_1+(1-r^2)^{3/2}V_2 \quad\cdots\text{①} \label{eq:3}
\end{align}
である．

\begin{figure}[htb]
\begin{center}
\begin{tikzpicture}[scale=2.2, >=stealth]

  % パラメータ設定 (例: r = 0.6 -> sqrt(1-r^2) = 0.8)
  \def\rVal{0.6}
  \def\rTwo{0.8}

  % ========================================================
  % 1. 左図: 球欠 V1 (交面 x = r, 積分区間 [r, 1])
  % ========================================================
  \begin{scope}[shift={(0, 0)}]
    \node[above, font=\bfseries] at (0, 1.35) {球 $C_1$ 側の正規化（$V_1$）};

    % 球欠部分の斜線ハッチング (r <= x <= 1)
    \begin{scope}
      \clip (\rVal, -1.0) rectangle (1.1, 1.0);
      \fill[pattern=north east lines, pattern color=blue!70] (0,0) circle (1.0);
    \end{scope}

    % 座標軸
    \draw[->, thick] (-1.3, 0) -- (1.4, 0) node[right] {$x$};
    \draw[->, thick] (0, -1.3) -- (0, 1.4) node[above] {$y$};
    \node[below left, font=\small] at (0,0) {$O$};

    % 単位円 x^2 + y^2 = 1
    \draw[ultra thick, blue!80!black] (0,0) circle (1.0);
    \node[blue!80!black, font=\small, above left] at (-0.6, 0.7) {$x^2+y^2=1$};

    % 切り口 x = r (破線)
    \pgfmathsetmacro{\yCutOne}{sqrt(1 - \rVal*\rVal)}
    \draw[thick, dashed, red!80!black] (\rVal, -\yCutOne) -- (\rVal, \yCutOne);
    \draw[dashed, gray] (\rVal, -\yCutOne) -- (\rVal, -1.1);

    % 目盛りとラベル
    \fill[red!80!black] (\rVal, 0) circle (1.2pt) node[below left, font=\small] {$r$};
    \fill[black] (1, 0) circle (1.0pt) node[below right, font=\small] {$1$};
    \fill[black] (0, 1) circle (1.0pt) node[above left, font=\small] {$1$};

    % 球欠ラベル
    \node[blue!80!black, font=\large] at (0.8, 0.3) {$V_1$};
  \end{scope}

  % ========================================================
  % 2. 右図: 球欠 V2 (交面 x = sqrt(1-r^2), 積分区間 [sqrt(1-r^2), 1])
  % ========================================================
  \begin{scope}[shift={(3.4, 0)}]
    \node[above, font=\bfseries] at (0, 1.35) {球 $C_2$ 側の正規化（$V_2$）};

    % 球欠部分の斜線ハッチング (sqrt(1-r^2) <= x <= 1)
    \begin{scope}
      \clip (\rTwo, -1.0) rectangle (1.1, 1.0);
      \fill[pattern=north east lines, pattern color=blue!70] (0,0) circle (1.0);
    \end{scope}

    % 座標軸
    \draw[->, thick] (-1.3, 0) -- (1.4, 0) node[right] {$x$};
    \draw[->, thick] (0, -1.3) -- (0, 1.4) node[above] {$y$};
    \node[below left, font=\small] at (0,0) {$O$};

    % 単位円 x^2 + y^2 = 1
    \draw[ultra thick, blue!80!black] (0,0) circle (1.0);
    \node[blue!80!black, font=\small, above left] at (-0.6, 0.7) {$x^2+y^2=1$};

    % 切り口 x = sqrt(1-r^2) (破線)
    \pgfmathsetmacro{\yCutTwo}{sqrt(1 - \rTwo*\rTwo)}
    \draw[thick, dashed, red!80!black] (\rTwo, -\yCutTwo) -- (\rTwo, \yCutTwo);
    \draw[dashed, gray] (\rTwo, -\yCutTwo) -- (\rTwo, -1.1);

    % 目盛りとラベル
    \fill[red!80!black] (\rTwo, 0) circle (1.2pt) node[below left, font=\small] {$\sqrt{1-r^2}$};
    \fill[black] (1, 0) circle (1.0pt) node[below right, font=\small] {$1$};
    \fill[black] (0, 1) circle (1.0pt) node[above left, font=\small] {$1$};

    % 球欠ラベル
    \node[blue!80!black, font=\large] at (0.9, 0.3) {$V_2$};
  \end{scope}

\end{tikzpicture}
\end{center}
\caption{単位球に正規化したときの球欠 $V_1$（区間 $[r, 1]$）と $V_2$（区間 $[\sqrt{1-r^2}, 1]$）の比較}
\end{figure}

球$C_1$の交面$x=r^2$は正規化すると$x'=r^2/r=r$に対応するので
\begin{align}
 \dfrac{V_1}{\pi}
  &=\int_r^1(1-x^2)\,dx \\
  &=\left[x-\dfrac{x^3}{3}\right]_r^1 \\
  &=\dfrac{2}{3}-r+\dfrac{1}{3}r^3 \label{eq:4}
\end{align}
である．

同様に球$C_2$の交面$x=r^2$は中心原点で正規化すると$x'=\dfrac{1-r^2}{\sqrt{1-r^2}}=\sqrt{1-r^2}$に対応するので\cref{eq:4}と同じ形で
\begin{align}
\dfrac{V_2}{\pi}
  &=\dfrac{2}{3}-\sqrt{1-r^2}+\dfrac{1}{3}\sqrt{1-r^2}^3 \label{eq:5}
\end{align}
である．

\cref{eq:4,eq:5}を\cref{eq:3}に代入して
\begin{align}
 \dfrac{V(r)}{\pi}
  &=r^3\left(\dfrac{2}{3}-r+\dfrac{1}{3}r^3\right)+\sqrt{1-r^2}^3\left(\dfrac{2}{3}-\sqrt{1-r^2}+\dfrac{1}{3}\sqrt{1-r^2}^3\right) \\
  &=\dfrac{2}{3}r^3-r^4+\dfrac{1}{3}r^6+\dfrac{2}{3}\sqrt{1-r^2}^3-(1-r^2)^2+\dfrac{1}{3}(1-r^2)^3 \\
\therefore
 V(r)&=\left(-r^4+\dfrac{2}{3}r^3+r^2-\dfrac{2}{3}+\dfrac{2}{3}(1-r^2)^{3/2}\right)\pi  \label{eq:6}
\end{align}
である．

\paragraph{(2)}

$V(r)$の挙動を調べるため一階微分を計算すると
\begin{align}
\begin{split}
\dfrac{V'(r)}{\pi}
  &=-4r^3+2r^2+2r+\dfrac{2}{3}\cdot\dfrac{3}{2}(1-r^2)^{1/2}\cdot(-2r) \\
  &=-4r^3+2r^2+2r-2r(1-r^2)^{1/2} \\
  &=-2r\left[2r^2-r-1+(1-r^2)^{1/2}\right]
\end{split}
\end{align}
である．
$0<r<1$では$-2r<0$だから，$f(r):=2r^2-r-1+\sqrt{1-r^2}$とおくと，$V'(r)$の符号は$-f(r)$の符号に等しい．この符号について
\begin{align}
 f(r)\ge0 \iff \sqrt{1-r^2} \ge 1+r-2r^2 
\end{align}
この右辺は$0<r<1$で正だから両辺二乗してよく，整理を続けて
\begin{align}
\begin{split}
 1-r^2& \ge(1+r-2r^2)^2\\
      &=4r^4-4r^2(r+1)+(r+1)^2 \\
      &=4r^4-4r^3-3r^2+2r+1 \\
\iff
 & 4r^4-4r^3-2r^2+2r\le 0 \\
\iff
 & 2r(2r^3-2r^2-r+1)\le0 \\
\iff
 &2r(r-1)(2r^2-1)\le 0
\end{split}
\end{align}
を得る．従って$V(r)$の増減表は以下のようになる．

\begin{table}[htb]
  \centering
\caption{$V(r)$の増減表}
\begin{tabular}{c|ccccc}
$r$ & $0$ & & $1/\sqrt{2}$ & & $1$ \\
\hline
$f$ & & $-$ & $0$ & $+$ & \\
\hline
$V'$ & & $+$ & $0$ & $-$ & \\
\hline
$V$ & & $\nearrow$ & & $\searrow$ &
\end{tabular}
\end{table}

よって$V(r)$は$r=1/\sqrt{2}=\sqrt{2}/2$で最大値をとる．\cref{eq:6}に代入してこの時の最大値は
$r^2=1/2,\ r^3=\sqrt{2}/4,\ r^4=1/4,\ 1-r^2=1/2,\ (1-r^2)^{3/2}=\sqrt{2}/4$だから
\begin{align}
\dfrac{V(1/\sqrt{2})}{\pi}
  &=-\dfrac{1}{4}+\dfrac{2}{3}\cdot\dfrac{\sqrt{2}}{4}+\dfrac{1}{2}-\dfrac{2}{3}+\dfrac{2}{3}\cdot\dfrac{\sqrt{2}}{4} \\
  &=\left(-\dfrac{1}{4}+\dfrac{1}{2}-\dfrac{2}{3}\right)+\dfrac{4}{3}\cdot\dfrac{\sqrt{2}}{4} \\
  &=-\dfrac{5}{12}+\dfrac{\sqrt{2}}{3} \\
\therefore
 \max V&= V\left(\dfrac{\sqrt{2}}{2}\right) = \left(\dfrac{\sqrt{2}}{3}-\dfrac{5}{12}\right)\pi
\end{align}
である．

\end{document}
